\section{Incidents and Recovery}\label{sec:incidents}

A successful Lean build establishes that the declarations in the checked
dependency graph are well typed. The four cases below concern boundaries it
does not reach: whether a theorem states the intended mathematics, whether
a refactor preserved every public declaration, and whether the runtime can
complete the next build. The evidence combines repository records with six
agent sessions selected for their reconstruction value, neither randomly
nor exhaustively, so the cases document observed failures and
recoveries rather than failure rates. PR numbers refer to the leray-hopf
repository~\cite{lerayhopf2026repo}; the cases do not depend on them.
% CLAIM: CLM-004, CLM-008

\subsection{Statement correctness}\label{sec:inc001}

\emph{Case A. A false generalization detected before release.}
The limit passage takes place in an evolution triple
$V \hookrightarrow H \hookrightarrow V'$; in the classical theory, $H$ is
the divergence-free $L^2$ space and $V$ its $H^1$ subspace. A standard
lemma states that if $u \in L^2(0,T;V)$ and its weak time derivative
satisfies $u' \in L^2(0,T;V')$, then $u$ agrees almost everywhere with a
function continuous into $H$; classically this is what gives a weak
solution a value at every time. The exponents matter because the proof
integrates the pairing $t \mapsto \langle u'(t), u(t)\rangle$ over time,
which Cauchy--Schwarz makes integrable when both exponents equal~2. A
public lemma in the project, declared over an abstract triple, asserted
the conclusion for arbitrary exponents $p,q\geq 1$, with
$u \in L^p(0,T;V)$ and $u' \in L^q(0,T;V')$, although the only proof
route examined was the case $p=q=2$. The declaration
carried an explicit \texttt{sorry}, and a later proof attempt checked only
that no capstone theorem depended on it. The false general statement
remained public for about a month and was repeated in project
documentation. % CLAIM: CLM-008, INC-001

\emph{Detection and response.}
During release preparation, a statement reviewer tested the boundary case
$p=q=1$. Take $H=\ell^2$ with a weighted triple in which the $n$th basis
vector has $V$-norm $2^n$ and $V'$-norm $2^{-n}$, and let $u$ consist of
spikes of unit $H$-norm along successive basis vectors, placed on
disjoint intervals of width $4^{-n}$ that accumulate at $t=0$. The widths
$4^{-n}$ against the $V$-norms $2^n$ make $u \in L^1(0,T;V)$, and the
$V'$-norms $2^{-n}$ make $u' \in L^1(0,T;V')$, so the stated hypotheses
hold, but the unit peaks accumulating at a point admit no
almost everywhere equal representative continuous into $H$. The problem
was therefore the proposition, not a difficult proof. The repair restricted
the declaration to $p=q=2$ and moved it outside the public imports
(PR~\#162). A follow-up added a statement card recording the exact type,
the literature reference, and the boundary cases checked, together with a
regression check that fails if the unrestricted signature reappears or a
placeholder lacks a card (PR~\#170).
% CLAIM: CLM-008, INC-001

\emph{Lesson.}
A placeholder records only that a proof is absent; it provides no evidence
that the proposition is true. Before proof search, reviewers should test the
boundary cases admitted by a public statement, especially when its
quantifiers or exponents exceed those used in the motivating argument.
% CLAIM: CLM-008, INC-001

\emph{Case B. A missing hypothesis found before merge.}
On $\mathbb{T}^3$ the convection term enters the Galerkin equations
through the trilinear form $b(u,v,w) = \int (u\cdot\nabla v)\cdot w$, and
the energy inequality rests on its antisymmetry $b(u,v,w) = -b(u,w,v)$ for
divergence-free $u$. The Galerkin scheme truncates Fourier modes to a
finite box, and the truncated form $b_n$ is a finite sum over pairs of
modes in the box. A second placeholder asserted antisymmetry of $b_n$ for
arbitrary $v$ and $w$. The standard proof, however, reindexes the
inner sum by an involution of the modes, which stays inside the box only
when the Fourier coefficients of $v$ and $w$ vanish outside it, that is,
when both lie in the finite projection space $V_n$. Two proof revisions had
attempted the unrestricted statement before an independent reviewer
identified this missing closure condition. % CLAIM: CLM-008, INC-002

\emph{Detection and response.}
The reviewer tested inputs outside $V_n$. A numerical diagnostic gave a
discrepancy between $b_n(u,v,w)$ and $-b_n(u,w,v)$ of approximately $80.7$
without the projection hypotheses and approximately $10^{-13}$ after
imposing them. These values corroborated the structural diagnosis; they
were not a certified counterexample. The repair added the two hypotheses
$P_n v = v$ and $P_n w = w$, where $P_n$ is the projection onto $V_n$,
completed the Lean proof, and removed the placeholder before merge
(PR~\#27). % CLAIM: CLM-008, INC-002

\emph{Lesson.}
For identities derived from truncation or projection, review should test
inputs outside the invariant subspace. Case B shows a defect caught by that
review before merge, whereas the analogous review in Case A came only at
release preparation. % CLAIM: CLM-005, CLM-008, INC-001, INC-002

\subsection{Public declarations during refactoring}

The construction of the convection form on $\mathbb{T}^3$ and on
$\mathbb{R}^3$ shared a layer of tensor and gluing lemmas that had been
written twice, once per domain. A refactor replaced both copies with
instances of a single parametrized module and was required to preserve
every public theorem of each domain file, most as thin instantiations of
the shared one (PR~\#120). The refactored code built successfully, but one
public theorem on the $\mathbb{R}^3$ side, which had no downstream users,
was omitted from the list of theorems to reconnect and disappeared, so
neither the Lean build nor an earlier review of module structure revealed
the omission. % CLAIM: CLM-008, INC-004

\emph{Detection and response.}
A separate reviewer compared the public declaration inventory before and
after the refactor, statement by statement. The missing theorem was
restored as a thin wrapper around the shared implementation with its
original statement, and the reviewer repeated the comparison before merge.
% CLAIM: CLM-008, INC-004

\emph{Lesson.}
When a refactor promises interface preservation, the review target must be
the complete public declaration set, including declarations with no
current users. A build tests dependencies. An inventory comparison tests
preservation. % CLAIM: CLM-008, INC-004

\subsection{Resource availability and orchestration}\label{sec:inc005}

The VPS container had an effective memory limit of 3.42\,GiB and no swap.
Review jobs from the previous night left fourteen helper process chains
running after their Git worktrees were inactive, consuming about
1--1.5\,GiB. At the same time, the orchestrator misclassified a slow agent
as inactive and assigned a replacement to the same Git worktree. The
replacement did not receive the shared cache rule and started a full
mathlib build. Under this load, Lean builds repeatedly ended with
exit code 137 and no compiler diagnostic. % CLAIM: CLM-008, INC-005

\emph{Detection and response.}
Monitoring had covered the visible Lean process rather than memory
available to the entire container. Roughly seventy minutes after an earlier
status check had reported the container as healthy on that basis, the
author's next status request reached the session only after an eight-minute
delay, and was followed at once by the author's explicit instruction to
kill stale processes. Only then did a memory check report 606\,MiB
available, and inspection of each process working directory identified the
fourteen stale chains. A cumulative kill counter maintained by the Linux
control group (cgroup) subsystem, read during the response, was taken at
the time as confirmation that the out of memory (OOM) killer had terminated
the failed builds; an audit of host kernel logs later overturned that
reading, because every counted kill belonged to a host-level OOM event two
days earlier and no kernel OOM occurred on the incident day. The mechanism
that terminated the builds remains unidentified. Removing the stale
processes raised available memory to 2.2\,GiB. The worktree returned to the
original agent, the duplicate stopped, and the build completed under a
global lock. The three PRs blocked by the failed builds merged later that
day (PRs~\#174--\#176).
% CLAIM: CLM-008, INC-005

\emph{Lesson.}
An orchestrator must observe the complete process tree and available
memory, not only the foreground prover. After the incident, helper-process
cleanup, cache reuse, build serialization, and a pre-build memory check
entered the standing instructions and dispatch prompts. Once
monitoring tracked memory
available to the whole container, it also detected a later peak caused by
one legitimate build, so measuring the resource directly is more general
than cleanup tailored to one cause. The
corrected diagnosis adds an evidence rule: a causal account assembled only
from observations inside the agent session can misattribute the mechanism,
so it must be checked against records outside the session.
% CLAIM: CLM-008, INC-005
